\clearpage
\SourceLeaf{95}%

\begin{center}
\large SCH\"ONES THEOREM DER WAHRSCHEINLICHKEITSRECHNUNG.
\end{center}

\noindent\hspace*{\parindent}Wenn [man setzt]:
\[
\int e^{itu}\varphi t\,dt = \psi u \mathord{.} \sqrt{2\pi},
\]
\noindent wobei das Integral von $t=-\infty$ bis $t=+\infty$ genommen wird, so gilt
\[
\int e^{-itu}\psi u\,du = \varphi t \mathord{.} \sqrt{2\pi},
\]
\noindent wobei das Integral ebenso genommen wird.

Die Begründung dieses Satzes ist in der folgenden Formel enthalten.
\[
\Sigma\,\varphi\!\left(t+\frac{2\pi k}{\omega}\right) = \frac{\omega}{2\pi}\Sigma\, e^{-ikt\omega}\int \varphi\theta\mathord{.} d\theta\mathord{.} e^{ik\theta\omega},
\]
\noindent Dabei setzt das Integralzeichen voraus, dass von $\theta=-\infty$ bis $\theta=+\infty$ integriert wird, und das $\Sigma$-Zeichen bezieht sich auf alle positiven und negativen ganzzahligen Werte von $k$. Man braucht darin nur $\omega$ unbegrenzt abnehmen zu lassen.

\Needspace{14\baselineskip}
\begin{center}
BEMERKUNGEN ZUM THEOREM DER WAHRSCHEINLICHKEITSRECHNUNG.
\end{center}

\noindent\hspace*{\parindent}Das vorliegende Theorem hat \textsc{Gauss} in den Einbanddeckel eines Buches mit dem Titel »Opuscula mathematica, 1799---1813« eingetragen. Eine genaue Angabe dazu, wann der Satz gefunden wurde, ist nicht überliefert.

Durch eine einfache Umrechnung kann man dem Theorem auch die folgende völlig gleichwertige Form geben.
\[
2\pi\varphi(t) = \int^{+\infty}_{-\infty}\!\left(du\int^{+\infty}_{-\infty}\!\cos u(t-x)\mathord{.}\varphi(x)\,dx\right).
\]
